\chapter{Conception et urbanisation du système}
\label{ch:conception}

Ce chapitre décrit l'architecture déployée, l'urbanisation MADI du SI, puis les diagrammes UML retenus. Le positionnement marché est traité au chapitre~\ref{ch:etat-art} ; les acteurs et exigences au chapitre~\ref{ch:analyse}.

\section{Urbanisation MADI}
\label{sec:madi}

LexHub est un \textbf{SI de cabinet}, pas « une appli IA ». L'approche \textbf{MADI} (Métier, Applications, Données, Infrastructure) relie le process métier aux applications, aux maîtres de donnée et à l'infrastructure. La grille AS-IS / écart / TO-BE du benchmarking marché (chapitre~\ref{ch:etat-art}) s'applique ici à l'\textbf{intérieur} du SI.

\begin{tcolorbox}[colback=royal!4,colframe=royal!55,title=Règle d'urbanisation]
Processus maître : \textbf{dossier $\rightarrow$ pièces $\rightarrow$ rédaction $\rightarrow$ livrable}. L'IA est un service partagé ; Postgres est la vérité métier ; le LLM n'est jamais maître des faits.
\end{tcolorbox}

\subsection{Couche M --- Métier (AS-IS)}
\label{sec:madi-m}

Processus de valeur : onboarding ; ouverture de dossier ; dépôt de pièces ; rédaction / dictée de plaidoirie (\textbf{cœur}) ; agenda ; équipe ; abonnement ; administration plateforme (côté LexHub).

\subsection{Couche A --- Applications (AS-IS)}
\label{sec:madi-a}

Un produit, plusieurs applications logiques, un socle de services (auth, fichiers, IA).

\begin{table}[H]
\centering
\caption{Blocs applicatifs LexHub (AS-IS)}
\label{tab:madi-apps}
\small
\begin{tabularx}{\textwidth}{@{}>{\bfseries}p{3.1cm}X>{\ttfamily\small}p{3.6cm}@{}}
\toprule
Bloc & Rôle & Route / API typique \\
\midrule
Auth / compte & Login, OTP, invitation & \texttt{/login}, \texttt{auth.py} \\
Dossiers & Affaire, parties & \texttt{/cases}, \texttt{cases.py} \\
Documents & Pièces du dossier & \texttt{/documents}, \texttt{documents.py} \\
Rédaction / plaidoirie & Templates, dictée, révision & \texttt{/redaction}, \texttt{dictation.py} \\
Agenda & Planning & \texttt{/agenda} \\
Équipe & Rôles cabinet & \texttt{/team} \\
Billing & Plans, checkout & \texttt{/checkout}, \texttt{billing.py} \\
Admin plateforme & Cabinets, leads & \texttt{/admin} \\
IA (service partagé) & STT, LLM, OCR, synthèse & Workers + APIs IA \\
\bottomrule
\end{tabularx}
\end{table}

L'IA n'est pas une application métier : c'est un \textbf{catalogue de services} (transcrire, corriger, composer, extraire) appelé par les applis métier.

\subsection{Couche D --- Données (AS-IS)}
\label{sec:madi-d}

\begin{table}[H]
\centering
\caption{Maîtres de donnée LexHub}
\label{tab:madi-data}
\small
\begin{tabularx}{\textwidth}{@{}>{\bfseries}p{3.4cm}Xp{4.2cm}@{}}
\toprule
Donnée & Maître & Pas maître \\
\midrule
Cabinet (tenant) & PostgreSQL \texttt{cabinet} & Frontend, IA \\
Utilisateur / rôle & PostgreSQL + JWT & MinIO \\
Dossier / parties & PostgreSQL & Ce que le LLM « croit » \\
Document (métadonnées) & PostgreSQL & Fichier brut \\
Fichier binaire & MinIO (chiffré) & Postgres (id seulement) \\
Index de recherche & ChromaDB & Index seulement \\
Tâches en cours & Redis & Volatile \\
Transcription / plaidoirie & Postgres (+ fichier) & Le LLM n'invente pas les faits \\
\bottomrule
\end{tabularx}
\end{table}

Isolation multi-tenant par \texttt{cabinet\_id}. Règle : le LLM propose ; l'avocat valide ; Postgres enregistre.

\subsection{Couche I --- Infrastructure (AS-IS)}
\label{sec:madi-i}

\begin{table}[H]
\centering
\caption{Infrastructure LexHub (rôle urbain)}
\label{tab:madi-infra}
\small
\begin{tabularx}{\textwidth}{@{}>{\bfseries}p{3.6cm}X@{}}
\toprule
Brique & Rôle urbain \\
\midrule
Frontend Next.js & Guichet visible \\
Backend FastAPI & Règles et API \\
PostgreSQL 16 & Cadastre (données métier) \\
MinIO & Entrepôt de fichiers \\
Redis + Celery & Files (OCR, indexation, tâches) \\
Docker Compose (VPS) & Terrain d'exécution \\
OpenAI (STT + LLM) & Partenaire externe, pas le cadastre \\
SMTP / Resend & Courrier transactionnel \\
HTTPS / JWT / chiffrement fichiers & Sécurité \\
\bottomrule
\end{tabularx}
\end{table}

Le cloud IA est \textbf{extérieur} à la ville : on l'utilise ; on ne lui confie pas le cadastre. Si le partenaire est indisponible, dossiers et pièces doivent encore exister.

\subsection{Plan de ville AS-IS}
\label{sec:madi-plan}

\begin{figure}[H]
\centering
\begin{tikzpicture}[
  layer/.style={draw=royal!55,fill=royal!5,rounded corners=3pt,align=left,font=\small,inner sep=6pt},
  arr/.style={-{Stealth},thick,royal!70}
]
  \node[layer,minimum width=14.2cm] (m) at (0,4.2) {%
    \textbf{M --- Métier}\quad Dossier $\rightarrow$ Pièces $\rightarrow$ Plaidoirie $\rightarrow$ Équipe $\rightarrow$ Abonnement};
  \node[layer,minimum width=14.2cm] (a) at (0,2.7) {%
    \textbf{A --- Applis}\quad Cases \quad Documents \quad Rédaction/Dictée \quad Billing \quad Admin};
  \node[layer,minimum width=14.2cm] (d) at (0,1.2) {%
    \textbf{D --- Données}\quad Postgres (vérité) \quad MinIO (fichiers) \quad Chroma (index) \quad Redis (files)};
  \node[layer,minimum width=14.2cm] (i) at (0,-0.3) {%
    \textbf{I --- Infra}\quad Docker/VPS \quad OpenAI (externe) \quad HTTPS/JWT};
  \draw[arr] (m) -- (a);
  \draw[arr] (a) -- (d);
  \draw[arr] (d) -- (i);
\end{tikzpicture}
\caption{Cartographie MADI LexHub (AS-IS)}
\label{fig:madi-plan}
\end{figure}

\subsection{Cible TO-BE, roadmap et gouvernance}
\label{sec:madi-tobe}

Urbaniser $\neq$ « tout casser en microservices ». La cible, ce sont des \textbf{règles} : un dossier = un maître (\texttt{case\_id}) ; l'IA = service, pas silo ; LLM $\neq$ cadastre ; multi-tenant partout (\texttt{cabinet\_id}) ; fichier $\neq$ métadonnée (MinIO / Postgres).

Architecture logique cible : \textbf{socle} (identité, tenant, droits, billing) ; \textbf{cœur cabinet} (dossiers, documents, agenda, équipe) ; \textbf{cœur production} (templates, dictée, plaidoirie, export) ; \textbf{socle IA} (STT, OCR, RAG, révision).

\begin{table}[H]
\centering
\caption{Roadmap d'urbanisation LexHub}
\label{tab:madi-waves}
\small
\begin{tabularx}{\textwidth}{@{}c>{\bfseries}p{3.2cm}XX@{}}
\toprule
Vague & Quoi & Valeur métier & Urbanisation \\
\midrule
1 & Dossier + pièces + droits & Le cabinet travaille dedans & Cadastre + routes \\
2 & Plaidoirie / dictée liée au dossier & Cœur de valeur & Service métier exposé \\
3 & Référentiels (parties, templates, corpus) & Moins de ressaisie & Couche Données propre \\
4 & Catalogue d'API interne stable & Mobile, partenaires & Réutilisation \\
5 & Billing / admin / observabilité & Entreprise & Gouvernance + mesure \\
\bottomrule
\end{tabularx}
\end{table}

À la date du rapport, LexHub se situe entre les \textbf{vagues 2 et 3}. Le travail d'urbaniste consiste surtout à \textbf{empêcher la dégradation} (nouveaux écrans ou features IA sans maître de donnée).

\begin{table}[H]
\centering
\caption{Gouvernance d'urbanisation --- questions utiles}
\label{tab:madi-gov}
\small
\begin{tabularx}{\textwidth}{@{}>{\bfseries}p{3.4cm}X@{}}
\toprule
Rôle & Question \\
\midrule
Métier & Est-ce que ça aide à produire un acte \textbf{dans un dossier} ? \\
Architecte & Ça casse le maître de donnée ? Ça duplique un service ? \\
Ops & C'est déployable, sauvegardé, isolé par cabinet ? \\
\bottomrule
\end{tabularx}
\end{table}

\section{Architecture générale}
\label{sec:archi}

LexHub suit une architecture en couches, classique d'un SaaS web. L'\textbf{interface} (application Next.js dans le navigateur) n'établit \textbf{aucune connexion} vers PostgreSQL, Redis ou le stockage de fichiers. Toute lecture ou écriture transite par l'\textbf{API HTTP FastAPI} : le client envoie une requête authentifiée (jeton JWT), le serveur vérifie le rôle et le \texttt{cabinet\_id}, puis accède aux données. C'est cette API --- le backend LexHub --- que désigne la formule « tout transite par l'API ».

La figure~\ref{fig:archi} présente les blocs de production.

\begin{figure}[H]
\centering
\begin{tikzpicture}[
  box/.style={draw=royal!65,fill=royal!6,rounded corners=4pt,align=center,minimum width=2.55cm,minimum height=1.12cm,font=\small},
  ext/.style={draw=emerald!70,fill=emerald!8,rounded corners=4pt,align=center,minimum width=2.45cm,minimum height=1.12cm,font=\small}
]
  \node[box] (fe) at (0,0) {Navigateur\\Next.js};
  \node[box] (caddy) at (3.5,0) {Caddy\\TLS};
  \node[box] (api) at (7.0,0) {API FastAPI};
  \node[box] (pg) at (10.5,0) {PostgreSQL};
  \node[box] (stt) at (0,-2.15) {STT\\Web Speech / Whisper};
  \node[ext] (ocr) at (3.5,-2.15) {OCR cloud};
  \node[box] (rd) at (7.0,-2.15) {Redis\\tâches + temps réel};
  \node[box] (io) at (10.5,-2.15) {Fichiers\\(MinIO)};
  \node[ext] (llm) at (3.5,-4.2) {LLM\\rédaction / révision};
  \node[box] (cel) at (7.0,-4.2) {Celery\\workers};
  \node[ext] (mail) at (10.5,-4.2) {E-mail\\transactionnel};
  \draw[-{Stealth},thick,royal] (fe) -- (caddy);
  \draw[-{Stealth},thick,royal] (caddy) -- (api);
  \draw[-{Stealth},thick,royal] (api) -- (pg);
  \draw[-{Stealth},thick,royal] (api) -- (stt);
  \draw[-{Stealth},thick,royal] (api) -- (ocr);
  \draw[-{Stealth},thick,royal] (api) -- (rd);
  \draw[-{Stealth},thick,royal] (api) -- (io);
  \draw[-{Stealth},thick,royal] (api) -- (llm);
  \draw[-{Stealth},thick,royal] (api) -- (cel);
  \draw[-{Stealth},thick,royal] (api) -- (mail);
\end{tikzpicture}
\caption{Architecture déployée --- vue logique (\texttt{lexhub.space})}
\label{fig:archi}
\end{figure}

\paragraph{Rôle de chaque bloc.}
Caddy termine TLS et répartit vers le frontend et l'API. FastAPI orchestre métier, droits et appels sortants. PostgreSQL conserve le modèle relationnel. Redis sert les files Celery et le canal temps réel. MinIO (stockage objet) conserve pièces et exports. La dictée et l'OCR passent par des API cloud ; le modèle de langage structure la plaidoirie sans inventer le fond. Les e-mails (OTP, invitation, rappels) sont transactionnels. La page de connexion s'appuie, après trois échecs, sur Cloudflare Turnstile (vérification du jeton par l'API, pas par le navigateur seul).

\section{Données, stockage et isolation}
\label{sec:bdd}

Une application de cabinet n'existe pas sans un modèle de données explicite. LexHub s'appuie sur trois supports complémentaires.

\paragraph{PostgreSQL.}
La base relationnelle est la source de vérité : cabinets, utilisateurs, dossiers, brouillons de plaidoirie, invitations, abonnements, journaux. Chaque table métier porte un \texttt{cabinet\_id}. Les requêtes API filtrent systématiquement sur ce locataire : un avocat de Rabat ne peut pas lire un dossier d'un cabinet de Tanger, même s'il devine un identifiant. Les évolutions de schéma passent par Alembic (migrations versionnées).

\paragraph{Stockage objet.}
Les fichiers (pièces scannées, logos, exports) ne sont pas stockés dans PostgreSQL. Ils sont chiffrés côté application puis déposés dans un magasin d'objets (MinIO en production). La base ne conserve que les métadonnées et le chemin.

\paragraph{Redis.}
Files d'attente (Celery : OCR, e-mails, rappels) et notifications temps réel. Redis n'est pas la base métier.

\paragraph{Ce que le navigateur ne voit jamais.}
Le client web ne reçoit ni chaîne de connexion, ni mot de passe SQL. Visualiser la base n'est pas une fonction produit : c'est un geste d'exploitation. En production, PostgreSQL, Redis et MinIO n'écoutent pas sur l'Internet public ; leurs ports sont liés à \texttt{127.0.0.1} sur le VPS. L'admin plateforme ne « ouvre » pas les dossiers depuis l'UI.

\paragraph{Visualisation et audit sécurisé de la base de données (exploitation).}
La base de données relationnelle \emph{est} PostgreSQL. Les données hébergées étant couvertes par le secret professionnel le plus strict (article~30 de la loi n$^\circ$~28-08 régissant la profession d'avocat au Maroc), l'exposition directe d'une interface de gestion sur le Web public (\url{https://lexhub.space/_ops/db/}) a été \textbf{totalement coupée} au niveau du reverse-proxy Caddy (toute requête renvoie 404/308).

Pour l'audit et l'évaluation, la solution retenue en production est le gestionnaire web ultra-léger \textbf{Adminer}, conteneurisé sur le VPS :
\begin{itemize}[leftmargin=1.4em]
  \item \textbf{Pourquoi Adminer plutôt que pgAdmin :} pgAdmin exige une installation desktop lourde, des dépendances système complexes ou un conteneur Docker client local sur le poste de l'évaluateur. À l'inverse, Adminer ne nécessite strictement aucune installation ni conteneur local : il tourne directement dans le navigateur.
  \item \textbf{Tunnel SSH chiffré en 1 clic :} Un script automatisé \texttt{Ouvrir\_Adminer.bat} (ou \texttt{python scripts/open\_adminer\_tunnel.py}) monte une redirection de port chiffrée (\texttt{127.0.0.1:8088} vers le conteneur Adminer du VPS).
  \item \textbf{Authentification stricte de production :} Connexion sur \texttt{http://127.0.0.1:8088} avec le système PostgreSQL, serveur interne \texttt{postgres}, utilisateur \texttt{lexhub} et base \texttt{lexhub}.
  \item \textbf{Révocation immédiate :} La simple fermeture du terminal coupe instantanément la passerelle.
\end{itemize}

\paragraph{Sauvegardes planifiées.}
La continuité d'exploitation n'est pas laissée au hasard. Un \textbf{ordonnancement quotidien} (tâche planifiée sur le VPS) produit une copie de PostgreSQL et du magasin de fichiers, conservée plusieurs jours. Ce dispositif est distinct du déploiement continu (CI/CD) : l'un publie le code, l'autre protège les \emph{données} du cabinet. Les identifiants et le détail opératoire restent dans le dossier de passation interne, pas dans ce rapport.

\paragraph{Risques et mesures.}
Un port base ouvert sur $0.0.0.0$ serait un risque d'énumération. La composition de production force l'écoute locale. Les identifiants vivent dans les variables d'environnement du VPS, pas dans le dépôt. Un code de parrainage n'est \textbf{pas} un jeton d'accès : il n'authentifie personne et n'ouvre aucun dossier (section~\ref{sec:parrainage}).

\section{Pipeline de traitement de la plaidoirie}

Les cinq traitements de l'atelier (dossier et pièces ; capture vocale ; correction de forme ; rédaction assistée ; export Word) sont formalisés par le diagramme d'activité (figure~\ref{fig:activite}) et détaillés côté usage au chapitre~\ref{ch:realisation}.

\section{Diagrammes UML retenus}

Le rapport documente les \textbf{séquences principales} de LexHub : inscription OTP, login / Turnstile, dictée~$\rightarrow$~rédaction, validation / export, dépôt document OCR, et facturation --- sources PlantUML dans \texttt{diagrams/}. S'y ajoutent le diagramme de \textbf{cas d'utilisation} et la \textbf{matrice RBAC} (chapitre~\ref{ch:analyse}), un diagramme d'\textbf{activité} (atelier), un diagramme de \textbf{classes} (extrait persistant), un diagramme de \textbf{déploiement}, et la vue d'\textbf{architecture} logique (figure~\ref{fig:archi}).

\subsection{Séquence --- de la voix au texte corrigé}

\begin{figure}[H]
\centering
\begin{tikzpicture}[font=\scriptsize, x=2.2cm, y=-0.48cm]
  \foreach \x/\lab in {0/{Avocat}, 1/{Navigateur}, 2/{API\\FastAPI}, 3/{Whisper\\(STT)}, 4/{LLM}, 5/{PostgreSQL}} {
    \node[align=center,font=\small\bfseries] at (\x,0) {\lab};
    \draw[linegray,thick,dashed] (\x,0.65) -- (\x,13.0);
  }

  \draw[fill=royal!15,draw=royal] (-0.08,1.0) rectangle (0.08,12.4);
  \draw[fill=royal!15,draw=royal] (0.92,1.1) rectangle (1.08,12.0);
  \draw[fill=royal!15,draw=royal] (1.92,3.6) rectangle (2.08,11.4);
  \draw[fill=emerald!18,draw=emerald] (2.92,5.0) rectangle (3.08,6.6);
  \draw[fill=emerald!18,draw=emerald] (3.92,8.0) rectangle (4.08,9.4);
  \draw[fill=royal!15,draw=royal] (4.92,10.0) rectangle (5.08,11.0);

  \draw[-{Stealth},thick] (0.08,1.3) -- node[above,font=\tiny]{1. Micro (Chrome/Edge ou Firefox)} (0.92,1.3);
  \draw[dashed,-{Stealth}] (0.92,2.2) -- node[above,font=\tiny]{2a. Web Speech : texte live (Chrome/Edge)} (0.08,2.2);
  \draw[-{Stealth},thick] (0.92,3.6) -- node[above,font=\tiny]{2b. MediaRecorder $\rightarrow$ audio (Firefox / fichier)} (1.92,3.6);
  \draw[-{Stealth},thick] (0.08,4.4) -- node[above,font=\tiny]{3. Terminé + JWT (si Whisper)} (1.92,4.4);
  \draw[-{Stealth},thick] (2.08,5.2) -- node[above,font=\tiny]{4. Transcription Whisper} (2.92,5.2);
  \draw[dashed,-{Stealth}] (2.92,6.2) -- node[above,font=\tiny]{5. Texte brut (+ langue auto FR/AR)} (2.08,6.2);
  \draw[dashed,-{Stealth}] (1.92,7.0) -- node[above,font=\tiny]{6. Affichage transcription} (0.92,7.0);
  \draw[-{Stealth},thick] (0.08,8.0) -- node[above,font=\tiny]{7. Corriger la forme (JWT)} (1.92,8.0);
  \draw[-{Stealth},thick] (2.08,8.6) -- node[above,font=\tiny]{8. Prompt forme + dossier anonymisé} (3.92,8.6);
  \draw[dashed,-{Stealth}] (3.92,9.4) -- node[above,font=\tiny]{9. Fusḥā + surlignage STT} (2.08,9.4);
  \draw[-{Stealth},thick] (2.08,10.2) -- node[above,font=\tiny]{10. Persistance brouillon} (4.92,10.2);
  \draw[dashed,-{Stealth}] (4.92,10.8) -- node[above,font=\tiny]{11. OK} (2.08,10.8);
  \draw[dashed,-{Stealth}] (1.92,11.6) -- node[above,font=\tiny]{12. Texte corrigé à l'avocat} (0.92,11.6);
\end{tikzpicture}
\caption{Diagramme de séquence UML --- Dictée (Web Speech \emph{ou} Whisper) puis correction}
\label{fig:seq-stt}
\end{figure}

\noindent\textit{Note.} Les deux voies STT sont alternatives : Web Speech (aperçu instantané) \emph{ou} Whisper (Firefox / import audio). Ensuite : endpoints \texttt{/clarify} (QCM), \texttt{/compose} (rédaction ancrée) et \texttt{/revise} (révision assistée) --- voir \texttt{seq-dictation.puml}.

\subsection{Séquence --- validation et visibilité assistant}

\begin{figure}[H]
\centering
\begin{tikzpicture}[font=\scriptsize, x=2.35cm, y=-0.52cm]
  % Lignes de vie
  \foreach \x/\lab in {0/{Avocat}, 1/{API FastAPI}, 2/{PostgreSQL}, 3/{Redis\\Pub/Sub}, 4/{Assistant}} {
    \node[align=center,font=\small\bfseries] at (\x,0) {\lab};
    \draw[linegray,thick,dashed] (\x,0.7) -- (\x,11.8);
  }

  % Barres d'activation UML
  \draw[fill=royal!15,draw=royal] (-0.08,1.0) rectangle (0.08,4.0);
  \draw[fill=royal!15,draw=royal] (0.92,1.2) rectangle (1.08,6.8);
  \draw[fill=royal!15,draw=royal] (1.92,2.0) rectangle (2.08,3.6);
  \draw[fill=royal!15,draw=royal] (2.92,4.4) rectangle (3.08,6.0);
  \draw[fill=royal!15,draw=royal] (3.92,5.6) rectangle (4.08,11.2);
  \draw[fill=royal!15,draw=royal] (0.92,8.8) rectangle (1.08,10.6);

  % Messages
  \draw[-{Stealth},thick] (0.08,1.4) -- node[above,font=\tiny]{1. Valider le brouillon (assigner opt.)} (0.92,1.4);
  \draw[-{Stealth},thick] (1.08,2.2) -- node[above,font=\tiny]{2. UPDATE status = 'validated'} (1.92,2.2);
  \draw[dashed,-{Stealth}] (1.92,3.2) -- node[above,font=\tiny]{3. Statut persisté} (1.08,3.2);

  \draw[-{Stealth},thick] (1.08,4.6) -- node[above,font=\tiny]{4. Émission événement (draft:validated)} (2.92,4.6);
  \draw[-{Stealth},thick] (2.92,5.8) -- node[above,font=\tiny]{5. Notification temps réel WebSocket} (3.92,5.8);

  \draw[dashed,-{Stealth}] (0.92,6.6) -- node[above,font=\tiny]{6. Accusé de validation affiché} (0.08,6.6);

  \draw[-{Stealth},thick] (3.92,9.0) -- node[above,font=\tiny]{7. Demande export Word (.docx) \textbf{ou} PDF} (1.08,9.0);
  \draw[-{Stealth},thick] (1.08,9.6) -- node[above,font=\tiny]{8. Contrôle RBAC \& génération docx/PDF} (1.5,9.6) -- (1.5,10.2) -- (1.08,10.2);
  \draw[dashed,-{Stealth}] (1.08,10.6) -- node[above,font=\tiny]{9. Téléchargement document officiel} (3.92,10.6);
\end{tikzpicture}
\caption{Diagramme de séquence UML --- Validation et gouvernance collaborative}
\label{fig:seq-val}
\end{figure}

\noindent\textit{Note.} Après export réussi, le statut \texttt{finalized} est posé par un \texttt{PATCH} côté UI (\texttt{markFinalizedAfterExport}) --- distinct de la validation assistant.

\subsection{Séquence --- inscription OTP}

\begin{figure}[H]
\centering
\begin{tikzpicture}[font=\scriptsize, x=2.2cm, y=-0.48cm]
  \foreach \x/\lab in {0/{Futur\\admin}, 1/{Navigateur}, 2/{API}, 3/{E-mail}, 4/{PostgreSQL}} {
    \node[align=center,font=\small\bfseries] at (\x,0) {\lab};
    \draw[linegray,thick,dashed] (\x,0.65) -- (\x,11.2);
  }
  \draw[fill=royal!15,draw=royal] (-0.08,1.0) rectangle (0.08,10.4);
  \draw[fill=royal!15,draw=royal] (0.92,1.1) rectangle (1.08,10.0);
  \draw[fill=royal!15,draw=royal] (1.92,1.8) rectangle (2.08,9.6);
  \draw[fill=emerald!18,draw=emerald] (2.92,2.6) rectangle (3.08,4.0);
  \draw[fill=royal!15,draw=royal] (3.92,2.0) rectangle (4.08,8.8);

  \draw[-{Stealth},thick] (0.08,1.3) -- node[above,font=\tiny]{1. E-mail cabinet} (0.92,1.3);
  \draw[-{Stealth},thick] (0.92,2.0) -- node[above,font=\tiny]{2. POST /request-register-otp} (1.92,2.0);
  \draw[-{Stealth},thick] (2.08,2.8) -- node[above,font=\tiny]{3. Stocke OTP (hash, TTL)} (3.92,2.8);
  \draw[-{Stealth},thick] (2.08,3.6) -- node[above,font=\tiny]{4. Envoi code 6 chiffres} (2.92,3.6);
  \draw[dashed,-{Stealth}] (2.92,4.4) -- node[above,font=\tiny]{5. E-mail OTP} (0.08,4.4);
  \draw[-{Stealth},thick] (0.92,5.4) -- node[above,font=\tiny]{6. POST /verify-register-otp} (1.92,5.4);
  \draw[-{Stealth},thick] (2.08,6.0) -- node[above,font=\tiny]{7. Vérifie OTP} (3.92,6.0);
  \draw[-{Stealth},thick] (0.92,7.0) -- node[above,font=\tiny]{8. POST /register (cabinet + MDP)} (1.92,7.0);
  \draw[-{Stealth},thick] (2.08,7.8) -- node[above,font=\tiny]{9. Cabinet ACTIVE + ADMIN + essai} (3.92,7.8);
  \draw[dashed,-{Stealth}] (1.92,8.8) -- node[above,font=\tiny]{10. MessageResponse (pas de JWT)} (0.92,8.8);
  \draw[dashed,-{Stealth}] (0.92,9.8) -- node[above,font=\tiny]{11. Puis login $\rightarrow$ JWT} (0.08,9.8);
\end{tikzpicture}
\caption{Diagramme de séquence UML --- Inscription cabinet (OTP e-mail)}
\label{fig:seq-otp}
\end{figure}

\subsection{Séquence --- login + Turnstile}

\begin{figure}[H]
\centering
\begin{tikzpicture}[font=\scriptsize, x=2.2cm, y=-0.48cm]
  \foreach \x/\lab in {0/{User}, 1/{Navigateur}, 2/{API}, 3/{Turnstile}, 4/{PostgreSQL}} {
    \node[align=center,font=\small\bfseries] at (\x,0) {\lab};
    \draw[linegray,thick,dashed] (\x,0.65) -- (\x,10.6);
  }
  \draw[fill=royal!15,draw=royal] (-0.08,1.0) rectangle (0.08,9.8);
  \draw[fill=royal!15,draw=royal] (0.92,1.1) rectangle (1.08,9.4);
  \draw[fill=royal!15,draw=royal] (1.92,1.6) rectangle (2.08,9.0);
  \draw[fill=emerald!18,draw=emerald] (2.92,5.0) rectangle (3.08,7.2);
  \draw[fill=royal!15,draw=royal] (3.92,3.2) rectangle (4.08,8.6);

  \draw[-{Stealth},thick] (0.08,1.3) -- node[above,font=\tiny]{1. E-mail + mot de passe} (0.92,1.3);
  \draw[-{Stealth},thick] (0.92,2.0) -- node[above,font=\tiny]{2. GET /captcha-config} (1.92,2.0);
  \draw[dashed,-{Stealth}] (1.92,2.6) -- node[above,font=\tiny]{3. site\_key + seuil (3 échecs)} (0.92,2.6);
  \draw[-{Stealth},thick] (0.92,3.4) -- node[above,font=\tiny]{4. POST /login} (1.92,3.4);
  \draw[-{Stealth},thick] (2.08,4.2) -- node[above,font=\tiny]{5. Vérifie credentials ($<$ 3 échecs)} (3.92,4.2);
  \draw[-{Stealth},thick] (0.92,5.2) -- node[above,font=\tiny]{6. (si $\geq$ 3) challenge Turnstile} (2.92,5.2);
  \draw[dashed,-{Stealth}] (2.92,6.0) -- node[above,font=\tiny]{7. token} (0.92,6.0);
  \draw[-{Stealth},thick] (0.92,6.8) -- node[above,font=\tiny]{8. POST /login + turnstile\_token} (1.92,6.8);
  \draw[-{Stealth},thick] (2.08,7.4) -- node[above,font=\tiny]{9. verify\_turnstile} (2.92,7.4);
  \draw[-{Stealth},thick] (2.08,8.2) -- node[above,font=\tiny]{10. Credentials OK} (3.92,8.2);
  \draw[dashed,-{Stealth}] (1.92,9.0) -- node[above,font=\tiny]{11. JWT (access + refresh)} (0.92,9.0);
\end{tikzpicture}
\caption{Diagramme de séquence UML --- Connexion et Turnstile}
\label{fig:seq-login}
\end{figure}

\subsection{Séquence --- dépôt document OCR}

\begin{figure}[H]
\centering
\begin{tikzpicture}[font=\scriptsize, x=2.0cm, y=-0.48cm]
  \foreach \x/\lab in {0/{User}, 1/{Navigateur}, 2/{API}, 3/{MinIO}, 4/{Worker}, 5/{PostgreSQL}} {
    \node[align=center,font=\small\bfseries] at (\x,0) {\lab};
    \draw[linegray,thick,dashed] (\x,0.65) -- (\x,9.6);
  }
  \draw[fill=royal!15,draw=royal] (-0.08,1.0) rectangle (0.08,8.8);
  \draw[fill=royal!15,draw=royal] (0.92,1.1) rectangle (1.08,3.6);
  \draw[fill=royal!15,draw=royal] (1.92,1.6) rectangle (2.08,5.2);
  \draw[fill=royal!15,draw=royal] (2.92,2.2) rectangle (3.08,6.0);
  \draw[fill=emerald!18,draw=emerald] (3.92,4.0) rectangle (4.08,7.4);
  \draw[fill=royal!15,draw=royal] (4.92,2.8) rectangle (5.08,8.2);

  \draw[-{Stealth},thick] (0.08,1.3) -- node[above,font=\tiny]{1. Fichier (PDF, Word, image\ldots)} (0.92,1.3);
  \draw[-{Stealth},thick] (0.92,2.0) -- node[above,font=\tiny]{2. POST /upload (+ JWT)} (1.92,2.0);
  \draw[-{Stealth},thick] (2.08,2.6) -- node[above,font=\tiny]{3. Stocke objet} (2.92,2.6);
  \draw[-{Stealth},thick] (2.08,3.4) -- node[above,font=\tiny]{4. Métadonnées (extracted\_text=null)} (4.92,3.4);
  \draw[dashed,-{Stealth}] (1.92,4.2) -- node[above,font=\tiny]{5. DocumentResponse (en file)} (0.92,4.2);
  \draw[-{Stealth},thick] (2.08,5.0) -- node[above,font=\tiny]{6. Tâche d'extraction} (3.92,5.0);
  \draw[-{Stealth},thick] (3.92,5.8) -- node[above,font=\tiny]{7. Lit le fichier} (2.92,5.8);
  \draw[-{Stealth},thick] (4.08,6.6) -- node[above,font=\tiny]{8. OCR / parse $\rightarrow$ extracted\_text} (4.92,6.6);
  \draw[dashed,-{Stealth}] (4.92,7.6) -- node[above,font=\tiny]{9. Statut prêt} (1.92,7.6);
\end{tikzpicture}
\caption{Diagramme de séquence UML --- Dépôt de pièce et extraction OCR}
\label{fig:seq-ocr}
\end{figure}

\subsection{Séquence --- facturation}

\begin{figure}[H]
\centering
\begin{tikzpicture}[font=\scriptsize, x=2.4cm, y=-0.48cm]
  \foreach \x/\lab in {0/{Admin\\cabinet}, 1/{Navigateur}, 2/{API}, 3/{PostgreSQL}} {
    \node[align=center,font=\small\bfseries] at (\x,0) {\lab};
    \draw[linegray,thick,dashed] (\x,0.65) -- (\x,10.4);
  }
  \draw[fill=royal!15,draw=royal] (-0.08,1.0) rectangle (0.08,9.6);
  \draw[fill=royal!15,draw=royal] (0.92,1.1) rectangle (1.08,9.2);
  \draw[fill=royal!15,draw=royal] (1.92,1.6) rectangle (2.08,9.0);
  \draw[fill=royal!15,draw=royal] (2.92,2.2) rectangle (3.08,8.6);

  \draw[-{Stealth},thick] (0.08,1.3) -- node[above,font=\tiny]{1. Ouvre abonnement} (0.92,1.3);
  \draw[-{Stealth},thick] (0.92,2.0) -- node[above,font=\tiny]{2. GET /entitlements + /plans} (1.92,2.0);
  \draw[-{Stealth},thick] (2.08,2.8) -- node[above,font=\tiny]{3. Quotas / plans / packs} (2.92,2.8);
  \draw[dashed,-{Stealth}] (1.92,3.6) -- node[above,font=\tiny]{4. Offres affichées} (0.92,3.6);
  \draw[-{Stealth},thick] (0.92,4.6) -- node[above,font=\tiny]{5. POST /select-plan \emph{ou} credit-pack} (1.92,4.6);
  \draw[-{Stealth},thick] (2.08,5.4) -- node[above,font=\tiny]{6. Abonnement / crédits + facture} (2.92,5.4);
  \draw[-{Stealth},thick] (0.92,6.4) -- node[above,font=\tiny]{7. GET /history (filtre amount $> 0$)} (1.92,6.4);
  \draw[-{Stealth},thick] (2.08,7.2) -- node[above,font=\tiny]{8. Factures payées uniquement} (2.92,7.2);
  \draw[-{Stealth},thick] (0.92,8.2) -- node[above,font=\tiny]{9. GET /invoices/\{id\}/download} (1.92,8.2);
  \draw[dashed,-{Stealth}] (1.92,9.0) -- node[above,font=\tiny]{10. PDF facture} (0.92,9.0);
\end{tikzpicture}
\caption{Diagramme de séquence UML --- Abonnement, packs et factures ($> 0$~MAD)}
\label{fig:seq-billing}
\end{figure}

\subsection{Modèle de données relationnel}
\label{sec:modele-donnees}

La base PostgreSQL regroupe les tables métier autour du \textbf{cabinet} (locataire). Le tableau~\ref{tab:domaines-bdd} donne la cartographie par domaine ; la figure~\ref{fig:classes} ne retient que les \textbf{entités maîtresses} et leurs clés --- le détail des 29 tables figure en annexe technique / migrations Alembic, pas dans le schéma de conception.

\begin{table}[H]
\centering
\caption{Domaines de données LexHub (synthèse)}
\label{tab:domaines-bdd}
\small
\renewcommand{\arraystretch}{1.15}
\begin{tabular}{@{}>{\bfseries}p{3.4cm}p{5.2cm}p{5.0cm}@{}}
\toprule
Domaine & Tables principales & Rôle \\
\midrule
Tenant \& accès &
\texttt{cabinets}, \texttt{users}, \texttt{cabinet\_invites}, sessions / OTP &
Isolation multi-tenant, RBAC, invitations \\
\addlinespace
Dossier &
\texttt{cases}, \texttt{documents}, \texttt{agenda\_events} &
Affaire, pièces (métadonnées), planning \\
\addlinespace
Plaidoirie &
\texttt{plaidoirie\_drafts}, \texttt{plaidoirie\_revisions}, templates &
Brouillon $\rightarrow$ validé $\rightarrow$ finalisé \\
\addlinespace
Facturation &
plans, abonnements, packs, factures, quotas &
Essai, crédits, sièges \\
\addlinespace
Audit &
\texttt{audit\_logs}, archives &
Traçabilité (loi~28-08) \\
\bottomrule
\end{tabular}
\end{table}

La figure~\ref{fig:classes} est un \textbf{extrait relationnel} (ER) : une ligne = une FK. Les fichiers binaires restent dans MinIO ; Postgres ne conserve que le chemin et les métadonnées.

\begin{figure}[H]
\centering
\begin{tikzpicture}[
  font=\scriptsize,
  ent/.style={
    draw=royal!70, thick, rounded corners=3pt, fill=royal!4,
    align=left, inner sep=5pt, minimum width=3.55cm
  },
  lbl/.style={font=\tiny\bfseries, text=royal!80},
  ln/.style={-{Stealth}, thick, royal!55}
]
  % --- Ligne 1 : tenant ---
  \node[ent] (cab) at (0, 4.2) {%
    \textbf{Cabinet}\\
    \underline{id}\\
    nom, status, referral\_code};
  \node[ent] (usr) at (5.0, 4.2) {%
    \textbf{User}\\
    \underline{id} \quad cabinet\_id\\
    email, role};
  \node[ent] (plan) at (10.0, 4.2) {%
    \textbf{SubscriptionPlan}\\
    \underline{id}\\
    slug, prix, quotas};

  % --- Ligne 2 : métier ---
  \node[ent] (cas) at (0, 1.6) {%
    \textbf{Case}\\
    \underline{id} \quad cabinet\_id\\
    référence, matière, tribunal};
  \node[ent] (doc) at (5.0, 1.6) {%
    \textbf{Document}\\
    \underline{id} \quad case\_id\\
    storage\_path, extracted\_text};
  \node[ent] (sub) at (10.0, 1.6) {%
    \textbf{CabinetSubscription}\\
    \underline{id} \quad cabinet\_id\\
    plan\_id, status, crédits};

  % --- Ligne 3 : plaidoirie ---
  \node[ent] (drf) at (2.5, -1.0) {%
    \textbf{PlaidoirieDraft}\\
    \underline{id} \quad case\_id, user\_id\\
    status (draft/validated/finalized)\\
    assigned\_to\_user\_id};
  \node[ent] (rev) at (7.5, -1.0) {%
    \textbf{PlaidoirieRevision}\\
    \underline{id} \quad draft\_id\\
    body\_snapshot, created\_at};

  % Liens
  \draw[ln] (usr) -- node[lbl, above]{N:1} (cab);
  \draw[ln] (cas) -- node[lbl, left]{N:1} (cab);
  \draw[ln] (doc) -- node[lbl, above]{N:1} (cas);
  \draw[ln] (drf) -- node[lbl, left]{N:1} (cas);
  \draw[ln] (rev) -- node[lbl, above]{N:1} (drf);
  \draw[ln] (sub) -- node[lbl, right]{1:1} (cab);
  \draw[ln] (sub) -- node[lbl, right]{N:1} (plan);

  % Légende
  \node[anchor=west, font=\tiny, text=royal!70] at (-0.2, -2.5)
    {Règle : toute ligne métier porte \texttt{cabinet\_id}. MinIO = fichiers ; LLM $\neq$ source de vérité.};
\end{tikzpicture}
\caption{Schéma relationnel (extrait) --- entités maîtresses LexHub}
\label{fig:classes}
\end{figure}

\subsection{Activité --- atelier plaidoirie}
\label{sec:activite}

Conformément à la norme internationale OMG UML~2.5, le diagramme d'activité de l'atelier de plaidoirie (Figure~\ref{fig:activite}) modélise le flux d'exécution séquentiel et les embranchements décisionnels. Le point d'entrée est représenté par un \textbf{nœud initial} (disque plein noir), les étapes par des \textbf{actions} (rectangles à coins arrondis), les arbitrages par des \textbf{losanges de décision} dotés de conditions de garde explicites, et l'aboutissement par un \textbf{nœud final d'activité} (œil-de-bœuf constitué d'un cercle plein entouré d'un cercle extérieur).

\begin{figure}[H]
\centering
\begin{tikzpicture}[
  font=\small,
  initnode/.style={circle, fill=royal!90!black, inner sep=5.5pt},
  action/.style={rectangle, rounded corners=5pt, draw=royal, thick, fill=royal!5, align=center, minimum width=4.2cm, minimum height=0.75cm},
  decision/.style={diamond, draw=royal, thick, fill=gold!20, align=center, aspect=2.0, inner sep=2pt, font=\scriptsize},
  finalouter/.style={circle, draw=royal!90!black, thick, inner sep=6.0pt},
  finalinner/.style={circle, fill=royal!90!black, inner sep=3.5pt}
]
  % Nœud initial UML
  \node[initnode] (s) at (0,0) {};

  % Actions du workflow
  \node[action] (dos) at (0,-1.3) {Sélection dossier \& pièces (OCR)};
  \node[action] (d) at (0,-2.6) {Capture vocale (Web Speech API) / audio};
  \node[action] (c) at (0,-3.9) {Correction de forme \& arabe standard};
  
  % Décision 1 : Contrôle des contradictions
  \node[decision] (q1) at (0,-5.5) {Faits OK ?};
  \node[action] (r) at (4.5,-5.5) {Arbitrage avocat\\(feuille de clarification)};
  
  \node[action] (comp) at (0,-7.2) {Rédaction IA de l'acte \& révision};
  
  % Décision 2 : Validation déontologique
  \node[decision] (q2) at (0,-8.8) {Valider ?};
  \node[action] (wait) at (4.5,-8.8) {Maintien en brouillon\\(édition ultérieure)};
  
  \node[action] (exp) at (0,-10.5) {Génération \& export Word (.docx) / PDF};
  
  % Nœud final d'activité UML (œil-de-bœuf)
  \node[finalouter] (e_out) at (0,-11.9) {};
  \node[finalinner] (e_in) at (0,-11.9) {};

  % Transitions
  \draw[-{Stealth},thick] (s) -- (dos);
  \draw[-{Stealth},thick] (dos) -- (d);
  \draw[-{Stealth},thick] (d) -- (c);
  \draw[-{Stealth},thick] (c) -- (q1);
  
  \draw[-{Stealth},thick] (q1) -- node[above,font=\scriptsize]{\textbf{[Contradictions]}} (r);
  \draw[-{Stealth},thick] (r) |- (c);
  \draw[-{Stealth},thick] (q1) -- node[right,font=\scriptsize]{\textbf{[Faits cohérents]}} (comp);
  
  \draw[-{Stealth},thick] (comp) -- (q2);
  \draw[-{Stealth},thick] (q2) -- node[above,font=\scriptsize]{\textbf{[Brouillon]}} (wait);
  \draw[-{Stealth},thick] (q2) -- node[right,font=\scriptsize]{\textbf{[Validé]}} (exp);
  \draw[-{Stealth},thick] (exp) -- (e_out);
\end{tikzpicture}
\caption{Diagramme d'activité UML 2.5 --- Atelier de plaidoirie (norme stricte)}
\label{fig:activite}
\end{figure}

\subsection{Déploiement}
\label{sec:deploy-uml}

\begin{figure}[H]
\centering
\begin{tikzpicture}[
  node/.style={draw=royal, thick, rounded corners=3pt, align=center, fill=royal!5, font=\small, minimum width=2.7cm, minimum height=0.95cm},
  ext/.style={draw=emerald, thick, rounded corners=3pt, align=center, fill=emerald!8, font=\small, minimum width=2.5cm}
]
  \node[node] (user) at (0,2.35) {Navigateur};
  \node[node] (caddy) at (4.15,2.35) {Caddy\\TLS};
  \node[node] (fe) at (2.05,0.45) {Conteneur\\frontend};
  \node[node] (be) at (6.25,0.45) {Conteneur\\API FastAPI};
  \node[node] (wk) at (8.85,0.45) {Celery};
  \node[node] (pg) at (2.05,-1.55) {PostgreSQL};
  \node[node] (rd) at (5.05,-1.55) {Redis};
  \node[node] (io) at (8.05,-1.55) {MinIO};
  \node[ext] (stt) at (11.35,0.45) {Whisper /\\OCR cloud};
  \node[ext] (llm) at (11.35,-1.55) {LLM cloud};
  \draw[-{Stealth},thick] (user) -- (caddy);
  \draw[-{Stealth},thick] (caddy) -- (fe);
  \draw[-{Stealth},thick] (caddy) -- (be);
  \draw[-{Stealth},thick] (be) -- (pg);
  \draw[-{Stealth},thick] (be) -- (rd);
  \draw[-{Stealth},thick] (be) -- (io);
  \draw[-{Stealth},thick] (be) -- (wk);
  \draw[-{Stealth},thick] (be) -- (stt);
  \draw[-{Stealth},thick] (be) -- (llm);
\end{tikzpicture}
\caption{Déploiement production --- VPS, Caddy, Docker, services cloud}
\label{fig:deploy}
\end{figure}

Les images d'application (\texttt{frontend}, \texttt{backend}, \texttt{celery}) sont publiées par GitHub Actions vers GHCR, puis tirées sur le VPS. Caddy sert \texttt{lexhub.space} en HTTPS. Le détail opératoire du déploiement figure au chapitre~\ref{ch:tests}.

La lecture fonctionnelle du produit (atelier, équipe, pilotage, économie) est spécifiée au chapitre~\ref{ch:analyse} et illustrée au chapitre~\ref{ch:realisation} ; le modèle économique au chapitre~\ref{ch:economie}.
